% Options for packages loaded elsewhere
\PassOptionsToPackage{unicode}{hyperref}
\PassOptionsToPackage{hyphens}{url}
\documentclass[
]{article}
\usepackage{xcolor}
\usepackage{amsmath,amssymb}
\setcounter{secnumdepth}{-\maxdimen} % remove section numbering
\usepackage{iftex}
\ifPDFTeX
  \usepackage[T1]{fontenc}
  \usepackage[utf8]{inputenc}
  \usepackage{textcomp} % provide euro and other symbols
\else % if luatex or xetex
  \usepackage{unicode-math} % this also loads fontspec
  \defaultfontfeatures{Scale=MatchLowercase}
  \defaultfontfeatures[\rmfamily]{Ligatures=TeX,Scale=1}
\fi
\usepackage{lmodern}
\ifPDFTeX\else
  % xetex/luatex font selection
\fi
% Use upquote if available, for straight quotes in verbatim environments
\IfFileExists{upquote.sty}{\usepackage{upquote}}{}
\IfFileExists{microtype.sty}{% use microtype if available
  \usepackage[]{microtype}
  \UseMicrotypeSet[protrusion]{basicmath} % disable protrusion for tt fonts
}{}
\makeatletter
\@ifundefined{KOMAClassName}{% if non-KOMA class
  \IfFileExists{parskip.sty}{%
    \usepackage{parskip}
  }{% else
    \setlength{\parindent}{0pt}
    \setlength{\parskip}{6pt plus 2pt minus 1pt}}
}{% if KOMA class
  \KOMAoptions{parskip=half}}
\makeatother
\setlength{\emergencystretch}{3em} % prevent overfull lines
\providecommand{\tightlist}{%
  \setlength{\itemsep}{0pt}\setlength{\parskip}{0pt}}
\usepackage{bookmark}
\IfFileExists{xurl.sty}{\usepackage{xurl}}{} % add URL line breaks if available
\urlstyle{same}
\hypersetup{
  hidelinks,
  pdfcreator={LaTeX via pandoc}}

\author{}
\date{}

\begin{document}

\section{GCP Deployment Scenarios and Architecture Decisions
(One-Page)}\label{gcp-deployment-scenarios-and-architecture-decisions-one-page}

\subsection{Objective}\label{objective}

Deliver the same online boutique application across three deployment
scenarios on GCP while keeping checkout reliable, data durable, and
operations practical for different scale and team maturity levels.

\subsection{Terminology (Quick
Clarifications)}\label{terminology-quick-clarifications}

\begin{itemize}
\tightlist
\item
  Sidecar: an additional container running in the same runtime boundary
  as a primary service container. It provides supporting capabilities
  (for example, database connectivity proxy or co-located dependency)
  while communicating over localhost.
\item
  RBAC: Role-Based Access Control.
\item
  TLS: Transport Layer Security.
\item
  IAM: Identity and Access Management.
\item
  Pub/Sub: Publish/Subscribe messaging.
\end{itemize}

\subsection{Application Components}\label{application-components}

\begin{itemize}
\tightlist
\item
  User-facing: frontend
\item
  Core commerce services: productcatalogservice, cartservice (with
  redis-cart), recommendationservice, checkoutservice, paymentservice,
  shippingservice, currencyservice, emailservice, adservice
\item
  Data services: Cloud SQL for PostgreSQL (transactional orders),
  Firestore (inventory metadata/state)
\item
  Eventing: Google Cloud Pub/Sub topic order-notifications
\item
  Additional serverless components: Cloud Run service order-workflow,
  Cloud Function subscriber for notifications/workflows
\end{itemize}

\subsection{Scenario 1: VM Deployment (Compute Engine + Docker
Compose)}\label{scenario-1-vm-deployment-compute-engine-docker-compose}

\begin{itemize}
\tightlist
\item
  Runtime model: all microservices run as containers on one VM.
\item
  Networking: internal service-to-service calls on compose network names
  and ports.
\item
  Data path: checkoutservice writes to Cloud SQL (private IP or Cloud
  SQL proxy path).
\item
  Best fit: low complexity, lower traffic, fast debugging, direct host
  control.
\item
  Tradeoff: manual scaling, patching, and host lifecycle ownership.
\end{itemize}

\subsection{Scenario 2: Container Platform Deployment
(GKE)}\label{scenario-2-container-platform-deployment-gke}

\begin{itemize}
\tightlist
\item
  Runtime model: microservices as Kubernetes Deployments/Services.
\item
  Networking: Kubernetes service discovery and cluster networking.
\item
  Data path: checkoutservice uses Cloud SQL Auth Proxy sidecar for Cloud
  SQL access.
\item
  Best fit: production-grade orchestration, rolling updates,
  fine-grained controls.
\item
  Tradeoff: highest ops overhead (cluster lifecycle, Role-Based Access
  Control (RBAC), node and policy management).
\end{itemize}

\subsection{Scenario 3: Serverless Platform Deployment (Cloud
Run)}\label{scenario-3-serverless-platform-deployment-cloud-run}

\begin{itemize}
\tightlist
\item
  Runtime model: either single service (for order-workflow) or
  multi-container sidecar service for boutique components.
\item
  Networking:

  \begin{itemize}
  \tightlist
  \item
    Sidecar model: localhost ports between containers, Transport Layer
    Security (TLS) disabled for in-process hop.
  \item
    Split-service model: service.run.app:443 endpoints, Transport Layer
    Security (TLS) enabled.
  \end{itemize}
\item
  Data path: order-workflow on Cloud Run writes to Cloud SQL, updates
  Firestore, and publishes Google Cloud Pub/Sub events.
\item
  Best fit: variable traffic, rapid delivery, minimal infrastructure
  management.
\item
  Tradeoff: per-revision container limits, stricter runtime constraints,
  careful env/port/Transport Layer Security (TLS) alignment required.
\end{itemize}

\subsection{Why These Architecture Decisions Were
Made}\label{why-these-architecture-decisions-were-made}

\begin{enumerate}
\def\labelenumi{\arabic{enumi}.}
\tightlist
\item
  Cloud SQL (PostgreSQL) for order records
\end{enumerate}

\begin{itemize}
\tightlist
\item
  Reason: checkout and order history are strongly relational and require
  transactional integrity (order header + order items), SQL
  queryability, and predictable consistency.
\item
  Result: normalized tables (customer\_order, order\_item) and ACID
  semantics for financial events.
\end{itemize}

\begin{enumerate}
\def\labelenumi{\arabic{enumi}.}
\setcounter{enumi}{1}
\tightlist
\item
  Firestore for inventory metadata
\end{enumerate}

\begin{itemize}
\tightlist
\item
  Reason: inventory and product state updates are document-shaped and
  evolve quickly; schema flexibility and low-latency key access are
  valuable.
\item
  Result: simple product-centric reads/writes
  (inventory/\{product\_id\}) without relational migration friction.
\end{itemize}

\begin{enumerate}
\def\labelenumi{\arabic{enumi}.}
\setcounter{enumi}{2}
\tightlist
\item
  Why AlloyDB, MySQL, or SQL Server were not selected
\end{enumerate}

\begin{itemize}
\tightlist
\item
  AlloyDB (PostgreSQL-compatible): strong option, but this project did
  not require its premium performance profile or operational features at
  current scale; Cloud SQL for PostgreSQL provided sufficient
  transactional capability with lower complexity and cost for the
  assignment scope.
\item
  Cloud SQL for MySQL: technically feasible, but the existing schema and
  scripts align with PostgreSQL behavior and ecosystem; switching
  engines would add migration and compatibility effort without clear
  functional benefit.
\item
  Cloud SQL for SQL Server: not chosen due to higher licensing/cost
  considerations and unnecessary platform coupling for this
  microservices demo workload.
\item
  Result: Cloud SQL for PostgreSQL offered the best fit for
  compatibility, predictable operations, and rubric-focused
  implementation speed.
\end{itemize}

\begin{enumerate}
\def\labelenumi{\arabic{enumi}.}
\setcounter{enumi}{3}
\tightlist
\item
  Pub/Sub for decoupling post-checkout actions
\end{enumerate}

\begin{itemize}
\tightlist
\item
  Reason: order completion should not block on downstream notification
  or enrichment steps.
\item
  Result: asynchronous, resilient event-driven fanout with retry
  semantics.
\end{itemize}

\begin{enumerate}
\def\labelenumi{\arabic{enumi}.}
\setcounter{enumi}{4}
\tightlist
\item
  Additional serverless service (Cloud Function subscriber)
\end{enumerate}

\begin{itemize}
\tightlist
\item
  Reason: notification, audit, and downstream integration logic is
  event-driven and bursty; serverless is cost-efficient and
  operationally light.
\item
  Result: Cloud Function scales on demand from Pub/Sub events and keeps
  checkout latency stable.
\end{itemize}

\subsection{Pub/Sub Invocation Point and
Embedding}\label{pubsub-invocation-point-and-embedding}

\begin{itemize}
\tightlist
\item
  Is Pub/Sub invoked during checkout: yes, but after order placement has
  been accepted and persisted by order-workflow.
\item
  Which service publishes: order-workflow (Cloud Run) publishes
  ORDER\_CREATED to the order-notifications topic.
\item
  Is Pub/Sub embedded in checkoutservice: no. checkoutservice calls
  order-workflow; order-workflow owns publish logic so checkout can stay
  focused on transaction orchestration.
\item
  Why this matters: checkout latency remains stable because notification
  delivery is asynchronous.
\end{itemize}

\subsection{What ``Order Notification'' Means
Here}\label{what-order-notification-means-here}

\begin{itemize}
\tightlist
\item
  An order notification is an event payload (for example ORDER\_CREATED
  with order\_id, user\_email, totals, and items) emitted to Pub/Sub.
\item
  Consumers (such as a Cloud Function) process that event for
  non-blocking downstream actions: confirmation email, audit entry,
  analytics stream update, fraud rules, or fulfillment integration.
\item
  It is a technical event contract, not only an end-user email.
\end{itemize}

\subsection{Serverless Additional Service Interaction Across
Scenarios}\label{serverless-additional-service-interaction-across-scenarios}

Common pattern used by VM, GKE, and Cloud Run storefront variants: 1.
User checks out via frontend and checkoutservice. 2. checkoutservice
invokes order-workflow /order-events. 3. order-workflow persists order
to Cloud SQL and updates Firestore. 4. order-workflow publishes
ORDER\_CREATED to Pub/Sub. 5. Cloud Function subscriber consumes event
and performs email/SMS/CRM/audit actions.

Because this integration is event-driven and endpoint-based, the same
serverless function can support all three deployment scenarios without
changing core checkout behavior.

\subsection{Why Markdown (.md) Is Good for This
Deliverable}\label{why-markdown-.md-is-good-for-this-deliverable}

\begin{itemize}
\tightlist
\item
  Portable and lightweight plain text (easy to version in Git and review
  in pull requests).
\item
  Human-readable in raw form and well-rendered in GitHub/GitLab/VS Code.
\item
  Fast to update during architecture changes without proprietary tool
  lock-in.
\item
  Easy to export to Word/PDF when submission format requires it.
\end{itemize}

\subsection{Rubric Coverage
(Demonstrated)}\label{rubric-coverage-demonstrated}

\begin{itemize}
\tightlist
\item
  Architecture clarity: three deployment topologies and component
  boundaries are explicit.
\item
  Data design rationale: relational vs document storage justified by
  workload shape.
\item
  Security model: private DB connectivity, secret management, Identity
  and Access Management (IAM)-based service access.
\item
  Reliability: async eventing, retries, and reduced checkout critical
  path dependencies.
\item
  Scalability: VM (vertical/manual), GKE (horizontal/orchestrated),
  Cloud Run/Functions (autoscale).
\item
  Cost/operations: VM lowest abstraction, GKE highest control/overhead,
  serverless best for elastic demand.
\item
  Integration quality: one event contract supports all scenarios
  consistently.
\item
  Maintainability: shared business flow with environment-specific
  runtime choices.
\end{itemize}

\end{document}
